\documentclass[11pt]{article}
\usepackage[margin=1in]{geometry}
\usepackage{amsmath,amssymb,amsthm,mathtools}
\usepackage{booktabs}
\usepackage{microtype}
\usepackage[hidelinks]{hyperref}
\usepackage{enumitem}

\newtheorem{theorem}{Theorem}[section]
\newtheorem{proposition}[theorem]{Proposition}
\newtheorem{corollary}[theorem]{Corollary}
\newtheorem{lemma}[theorem]{Lemma}
\theoremstyle{definition}
\newtheorem{definition}[theorem]{Definition}
\theoremstyle{remark}
\newtheorem{remark}[theorem]{Remark}

\newcommand{\Q}{\mathbb Q}
\newcommand{\Z}{\mathbb Z}
\newcommand{\F}{\mathbb F}
\newcommand{\OO}{\mathcal O}
\newcommand{\tr}{\operatorname{tr}}
\newcommand{\Gal}{\operatorname{Gal}}
\newcommand{\disc}{\operatorname{disc}}

\title{\textbf{The Discriminant-$12$ Return}\\
\large From Centered Cayley Geometry to Pell, Cyclotomic, and Boolean Structure}
\author{Jeremy Ebert}
\date{2026}

\begin{document}
\maketitle

\begin{abstract}
Let $g\in SL_2$ have trace $\tau$, centered part
$H=g-\tfrac{\tau}{2}I$, and Cayley transform
$K=(g-I)(g+I)^{-1}$.  We prove the universal identities
\[
K=\frac{2}{\tau+2}H,
\qquad
KH=\frac{\tau-2}{2}I.
\]
Consequently, in the shear normalization $\tau=\sigma+2$, the
self-reciprocity condition $K=H^{-1}$ selects the unique value $\sigma=2$.
The period-$[1,2]$ primitive positive integral return is
\[
g_{12}=\begin{pmatrix}3&1\\2&1\end{pmatrix},
\]
whose characteristic discriminant is $12$ and whose eigenvalues are
$2\pm\sqrt3$.  We show that its centered generator realizes multiplication
by $\sqrt3$ on the ramified ideal $(2,\sqrt3-1)$, that the return preserves
the primitive form $2x^2-2xy-y^2$, and that this form is a Lorentz norm in
Pell coordinates.  The same centered generator admits the cyclotomic
completion $\mathcal Z=(H+iI)/2$, a primitive twelfth root whose integral
closure carries a four-state parity quotient.  Modulo $2$, the return becomes
factor exchange on a Boolean cell and its degenerate norm becomes XOR.
The argument is algebraic and does not depend on zeta zeros, numerical
certification, or the Riemann hypothesis.
\end{abstract}

\section{Introduction}

A hyperbolic matrix in $SL_2$ naturally produces two centered objects.  The
first is its traceless part
\[
H=g-\frac{\tr(g)}2I.
\]
The second is its Cayley transform
\[
K=(g-I)(g+I)^{-1}.
\]
Both lie on the one-dimensional traceless direction determined by $g$, but
their scalar relation depends on the trace.  Requiring them to be reciprocal
turns out to be rigid: after the normalization $\tr(g)=\sigma+2$, it forces
\[
\sigma=2.
\]

This elementary selection rule is the entrance to a longer arithmetic chain.
At the selected trace $4$, a primitive positive integral representative has
characteristic discriminant
\[
\Delta_g=\tr(g)^2-4=12
\]
and expanding eigenvalue $2+\sqrt3$.  Its centered part squares to $3I$; its
integral lattice is the ramified ideal above $2$ in $\Q(\sqrt3)$; its norm
form is an integral Lorentz form; and adjoining the scalar complex structure
produces a primitive twelfth root.  Reduction at the ramified prime returns
to a four-state Boolean cell, not as a translation but as factor exchange.

The purpose of this paper is to prove that chain in the order in which its
dependencies actually occur:
\[
\boxed{
\begin{gathered}
\text{centered Cayley reciprocity}
\Longrightarrow \sigma=2
\Longrightarrow \Delta_g=12
\Longrightarrow \Q(\sqrt3)\\
\Longrightarrow \text{Pell--Lorentz lattice}
\Longrightarrow \Q(\zeta_{12})
\Longrightarrow \text{ramified Boolean factor exchange}.
\end{gathered}}
\]

Three distinctions will be maintained throughout.  First,
$\Delta_g=\tr(g)^2-4$ is the characteristic discriminant of a matrix; it is
not automatically the discriminant of an arbitrarily chosen order.  Second,
the Lorentz dynamics is a discrete hyperbolic orbit, not a closed continuous
orbit.  Third, the Boolean translation group, the returned factor-exchange
involution, and the cyclotomic Galois group are separate symmetry layers.

The algebraic theorem proved here is independent of any analytic realization.
In particular, no identification with a Fredholm flow, no statement about
zeta zeros, and no implication for the Riemann hypothesis is used or claimed.

\section{Statement of the discriminant-$12$ return theorem}

\begin{theorem}[Discriminant-$12$ return]
\label{thm:main}
Let $g\in SL_2(\Q)$, let $\tau=\tr(g)$, and assume $g+I$ is invertible.
Define
\[
H=g-\frac\tau2I,
\qquad
K=(g-I)(g+I)^{-1},
\qquad
\tau=\sigma+2.
\]
Then:
\begin{enumerate}[label=(\roman*)]
\item $K=\dfrac{2}{\tau+2}H$ and
      $KH=\dfrac{\tau-2}{2}I$.
\item In the hyperbolic range, $K=H^{-1}$ if and only if $\sigma=2$.
\item At $\sigma=2$, the period-$[1,2]$ primitive positive integral
representative
\[
g_{12}=\begin{pmatrix}3&1\\2&1\end{pmatrix}
\]
has characteristic polynomial $X^2-4X+1$, characteristic discriminant $12$,
and eigenvalues $2\pm\sqrt3$.  Its centered part
\[
H_{12}=g_{12}-2I=\begin{pmatrix}1&1\\2&-1\end{pmatrix}
\]
satisfies $H_{12}^2=3I$.
\item On the lattice $\mathfrak p_2=(2,\sqrt3-1)$, written in the ordered
basis $(2,\sqrt3-1)$, the operators $H_{12}$ and $g_{12}$ are multiplication
by $\sqrt3$ and $2+\sqrt3$, respectively.  The associated primitive form
\[
q_{12}(x,y)=2x^2-2xy-y^2
\]
is preserved by $g_{12}$ and obeys
\[
2q_{12}(x,y)=U^2-V^2,
\qquad
(U,V)=(2x-y,\sqrt3\,y).
\]
\item The element
\[
\mathcal Z=\frac{H_{12}+iI}{2}
\]
satisfies $\Phi_{12}(\mathcal Z)=0$ and generates
$\Q(\zeta_{12})$.  Moreover
\[
g_{12}=-i(I+\mathcal Z)(I-\mathcal Z)^{-1}.
\]
The maximal order is
\[
\OO_{\Q(\zeta_{12})}
=\left\{
\frac{A+B\sqrt3+i(C+D\sqrt3)}2:
A\equiv D,\ B\equiv C\pmod2
\right\},
\]
and
\[
\OO_{\Q(\zeta_{12})}/\Z[\sqrt3,i]\cong(\F_2)^2.
\]
\item Modulo $2$,
\[
\overline g_{12}(x,y)=(x+y,y),
\qquad
\overline q_{12}(x,y)=y.
\]
Under $\Phi(x,y)=(\epsilon_1,\epsilon_2)=(x,x+y)$ these become
\[
\Phi\overline g_{12}\Phi^{-1}
(\epsilon_1,\epsilon_2)=(\epsilon_2,\epsilon_1),
\qquad
\overline q_{12}=\epsilon_1\oplus\epsilon_2.
\]
Thus the return is factor exchange and, together with Boolean translations,
generates $V_4^{\rm tr}\rtimes C_2\cong D_8$.
\end{enumerate}
\end{theorem}

The theorem separates into four proof packages: universal centered-Cayley
algebra, the integral Pell--Lorentz return, cyclotomic integral closure, and
the ramified Boolean reduction.  The first two are proved next.

\section{Universal centered-Cayley algebra}

The selection of the trace defect is a consequence of Cayley--Hamilton alone.

\begin{lemma}[Centered square]
\label{lem:centered-square}
For $g\in SL_2$ with trace $\tau$,
\[
\boxed{H^2=\left(\frac{\tau^2}{4}-1\right)I.}
\]
\end{lemma}

\begin{proof}
Cayley--Hamilton gives $g^2-\tau g+I=0$.  Therefore
\[
H^2
=\left(g-\frac\tau2I\right)^2
=g^2-\tau g+\frac{\tau^2}{4}I
=\left(\frac{\tau^2}{4}-1\right)I.
\]
\end{proof}

\begin{proposition}[Centered Cayley identity]
\label{prop:centered-cayley}
If $g+I$ is invertible, then
\[
\boxed{
K=\frac{2}{\tau+2}H,
\qquad
KH=\frac{\tau-2}{2}I.
}
\]
\end{proposition}

\begin{proof}
From Cayley--Hamilton,
\[
(g+I)\bigl((\tau+1)I-g\bigr)
=(\tau+2)I.
\]
Hence
\[
(g+I)^{-1}=\frac{(\tau+1)I-g}{\tau+2}.
\]
Multiplication gives
\begin{align*}
K
&=\frac{(g-I)((\tau+1)I-g)}{\tau+2}\\
&=\frac{(\tau+2)g-g^2-(\tau+1)I}{\tau+2}\\
&=\frac{2g-\tau I}{\tau+2}
=\frac{2}{\tau+2}H.
\end{align*}
Using Lemma~\ref{lem:centered-square},
\[
KH
=\frac{2}{\tau+2}H^2
=\frac{2}{\tau+2}\left(\frac{\tau^2}{4}-1\right)I
=\frac{\tau-2}{2}I.
\]
\end{proof}

\begin{corollary}[Unique self-reciprocal class]
\label{cor:self-reciprocal}
Let $\tau=\sigma+2$ and assume $g$ is hyperbolic.  Then
\[
\boxed{K=H^{-1}\quad\Longleftrightarrow\quad\sigma=2.}
\]
\end{corollary}

\begin{proof}
Hyperbolicity implies $H$ is invertible.  By
Proposition~\ref{prop:centered-cayley},
\[
K=H^{-1}
\quad\Longleftrightarrow\quad
KH=I
\quad\Longleftrightarrow\quad
\frac{\tau-2}{2}=1.
\]
Thus $\tau=4$, equivalently $\sigma=2$.
\end{proof}

\begin{remark}
The number $2$ is selected here before any integral representative, quadratic
field, Pell equation, or analytic flow has been introduced.  Later
appearances of $2+\sqrt3$ and discriminant $12$ are downstream consequences
of this trace selection.
\end{remark}

\section{The primitive integral return}

For the standard continued-fraction matrix
\[
A_a=\begin{pmatrix}a&1\\1&0\end{pmatrix},
\]
the period $[1,2]$ gives, at the selected trace $4$,
\[
g_{12}=A_1A_2=\begin{pmatrix}3&1\\2&1\end{pmatrix}.
\]
It is positive, primitive, and unimodular.  Directly,
\[
\det(g_{12})=1,
\qquad
\tr(g_{12})=4,
\]
so
\[
\chi_{g_{12}}(X)=X^2-4X+1,
\qquad
\Delta_{g_{12}}=4^2-4=12.
\]
The eigenvalues are $2\pm\sqrt3$.

\begin{proposition}[Real quadratic generator]
\label{prop:real-generator}
For
\[
H_{12}=g_{12}-2I
=\begin{pmatrix}1&1\\2&-1\end{pmatrix},
\]
one has
\[
H_{12}^2=3I,
\qquad
\Z[H_{12}]\cong\Z[\sqrt3]=\OO_{\Q(\sqrt3)}.
\]
\end{proposition}

\begin{proof}
Matrix multiplication gives $H_{12}^2=3I$.  Thus evaluation
$\sqrt3\mapsto H_{12}$ defines an injective homomorphism
$\Z[\sqrt3]\to M_2(\Z)$ with image $\Z[H_{12}]$.  Since
$3\equiv3\pmod4$, the maximal order of $\Q(\sqrt3)$ is
$\Z[\sqrt3]$.
\end{proof}

Let
\[
\mathfrak p_2=(2,\sqrt3-1).
\]
This is the unique prime ideal above $2$; indeed $2$ ramifies in the real
quadratic field of discriminant $12$.

\begin{proposition}[Ideal-lattice action]
\label{prop:ideal-action}
In the ordered basis
\[
e_1=2,
\qquad
e_2=\sqrt3-1
\]
of $\mathfrak p_2$, multiplication by $\sqrt3$ and $2+\sqrt3$ is represented
by $H_{12}$ and $g_{12}$, respectively.
\end{proposition}

\begin{proof}
The basis relations are
\[
\sqrt3e_1=e_1+2e_2,
\qquad
\sqrt3e_2=e_1-e_2.
\]
Therefore the columns of the multiplication matrix are $(1,2)^T$ and
$(1,-1)^T$, giving $H_{12}$.  Since multiplication by $2+\sqrt3$ is
$2I+H_{12}$, its matrix is $g_{12}$.
\end{proof}

\section{The Pell--Lorentz norm form}

The ideal norm becomes an integral indefinite binary quadratic form in the
basis above.  For
\[
\alpha=xe_1+ye_2=(2x-y)+y\sqrt3,
\]
one has
\[
N_{\Q(\sqrt3)/\Q}(\alpha)
=(2x-y)^2-3y^2
=2\bigl(2x^2-2xy-y^2\bigr).
\]
Define
\[
q_{12}(x,y)=2x^2-2xy-y^2.
\]
Its binary-form discriminant is
\[
(-2)^2-4(2)(-1)=12.
\]

\begin{proposition}[Integral Lorentz return]
\label{prop:lorentz-return}
The form $q_{12}$ is primitive and is preserved by $g_{12}$.  Under
\[
(U,V)=(2x-y,\sqrt3\,y),
\]
it satisfies
\[
\boxed{2q_{12}(x,y)=U^2-V^2.}
\]
Moreover $g_{12}$ acts as the discrete Lorentz boost associated with the unit
$2+\sqrt3$.
\end{proposition}

\begin{proof}
Primitivity follows from the coefficients $2,-2,-1$.  By
Proposition~\ref{prop:ideal-action}, $g_{12}$ is multiplication by the norm-one
unit $2+\sqrt3$, so it preserves the ideal norm and hence $q_{12}$.  The
displayed coordinate change is the norm calculation above.  In null
coordinates $U\pm V$, multiplication by $2+\sqrt3$ scales the two eigendirections
by $2\pm\sqrt3$.  Thus the iterates of $g_{12}$ form a discrete hyperbolic
Lorentz orbit on each level of $q_{12}$.
\end{proof}

\begin{remark}
The trace determines the recurrence law for the iterates, but not their
initial data.  Since $g_{12}$ has characteristic polynomial
$X^2-4X+1$, every scalar coordinate of $g_{12}^nv$ obeys
\[
a_{n+2}=4a_{n+1}-a_n,
\]
with initial values determined by $v$.
\end{remark}

\section{Roadmap for the arithmetic completion}

The remaining proof packages proceed without altering the hypotheses above.
First, adjoining the scalar complex structure gives
\[
\mathcal Z=\frac{H_{12}+iI}{2},
\qquad
\mathcal Z\overline{\mathcal Z}
=\frac{H_{12}^2+I}{4}=I.
\]
The half-element is therefore forced by relative-norm normalization and
satisfies the twelfth cyclotomic polynomial.  The integral closure introduces
two parity-gluing classes, yielding a quotient isomorphic as an additive group
to $(\F_2)^2$.

Second, reduction of the ramified ideal lattice modulo $2$ gives another
four-state space.  The return reduces to the transvection
\[
(x,y)\longmapsto(x+y,y),
\]
which becomes factor exchange after the Boolean relabeling
$(x,y)\mapsto(x,x+y)$.  Its reduced norm becomes XOR.  Equal cardinality does
not canonically identify this ramified quotient with the cyclotomic gluing
quotient, and neither should be identified with the cyclotomic Galois $V_4$.

These two packages complete parts (v) and (vi) of
Theorem~\ref{thm:main}.  They will be developed next from the integral parity
description rather than inferred from the shared number of states.

\section*{Scope note}

This draft establishes the theorem-first opening and the complete proofs of
the centered-Cayley and Pell--Lorentz packages.  A separate analytic companion
may realize the selected value through a certified Suzuki determinant
crossing.  That realization is not a hypothesis of the algebraic theorem and
is not used here.

\end{document}
